% All eight P76 conditions; generated only from the closed audited batch.
\section{推論履歴を保持する通常対照}
同じ既知４課題の８応答後から，公開履歴のみ（P）と返却推論も全て保持（R）を各１回比較した．
両条件は公開履歴・持続Python・公開名接続・予算を共通にし，過去と以後の推論の投入だけを変えた．
各16生成・120,000 token・480秒で，最初の生成はcacheを使わない．
共通前半は推論を再投入せず生成しており，最初からの方針比較ではない．
\begin{center}\begin{minipage}{\columnwidth}\centering
\small 推論保持比較の全８条件：成否・仕事token\par\smallskip
\begin{tabular}{lrr}\hline 課題 & P：公開履歴のみ & R：推論も保持\\\hline
写真 & 成功 156,754 & 成功 118,184\\
精算 & 未達 160,878 & 未達 164,322\\
返金 & 成功 85,720 & 成功 88,001\\
書出 & 未達 134,971$+?$ & 未達 100,425$+?$\\
\hline
成功／評価 & 2/4 & 2/4\\
既知token合計 & 538,323$+?$ & 470,932$+?$\\
\hline\end{tabular}
\end{minipage}\end{center}
Rの初回入力は課題順に291・189・1,116・361 token増えた．
全８予定条件のうち8条件を評価し，使用量不明は2生成だった．欠測を補完せず，採用判定に用いる確定費用差は求めない．
課題別のR−Pの仕事token差は，写真 -38,570／精算 +3,444／返金 +2,281／書出し 未確定だった．
全８評価・全使用量既知・少なくとも１対成功・成功喪失０・全体費用減少という局所採用条件は満たさなかった．
新規継続は43生成・既知588,495 token＋使用量不明2生成，継続306・復元102 APIだった．
推論器全稼働2,662.086秒に復元・評価も含め，重複加算しない．
保持は通常の比較条件であり，新規手法や一般的な非劣性の実証とはしない．表2・3の過去条件は変更しない．
